\appendix
\chapter{Domande delle lezioni}
\begin{multicols}{2}
\begin{enumerate}
    \item I sistemi biologici sono complessi 
    \item Il modello è una rappresentazione esatta della complessità biologica  
    \item Tutti i sistemi biologici hanno un comportamento non lineare che il modello può approssimare al caso lineare sotto certe ipotesi 
    \item Il corpo umano è composto da parti il cui comportamento può essere predetto dall'osservazione delle sue caratteristiche generali (olonomicità)
    \item Il feedback negativo è tipico del metabolismo del glucosio  Il controllo nell'organismo umano viene eseguito solo mediante enzimi 
    \item La ridondanza permette di mantenere la funzione in caso di danno  Per quantificare i processi biologici 3 necessario fare misure 
    \item Il modello rappresenta la realtà senza alcuna approssimazione Possibili tipologie di modello 
    \item La formulazione del modello dipende dall'obiettivo per cui viene definito
    \item Il modello interpretativo permette unicamente di descrivere il fenomeno
    \item Il modello permette di inferire misure che non possono essere realizzate in vivo  
    \item Nel processo di modellizzazione si può decidere di modellizzare i dati o il sistema stesso 
    \item Quali azioni permettono di approssimare la realtà in un modello  
    \item L'identificabilità del modello consiste nella formulazione della sua struttura  
    \item Il modello è valido in relazione agli obiettivi per cui è stato formulato  
    \item Il modello usato in modalità simulazione fornisce informazioni riguardo il comportamento del sistema
    \item Il modello di sistema tenta di rappresentare esplicitamente la fisiologia del sistema 
    \item Il modello di sistema può essere contemporaneamente distribuito e concentrato 
    \item Un modello è dinamico se le equazioni che lo descrivono contengono le derivate rispetto alla variabile spaziale 
    \item Un modello può essere contemporaneamente dinamico e non lineare 
    \item Nel modello stocastico le equazioni contengono elementi probabilistici  
    \item Il processo di identificazione permette di determinare completamente tutti i parametri definiti nella struttura del modello 
    \item L'ampiezza dei segnali di test non è un parametro critico nel processo di identificazione 
    \item I segnali di test transitori non eseguono nessun filtraggio del rumore 
    \item Applicare un segnale di test random equivale ad eseguire test armonici multipli contemporaneamente 
    \item   L'errore di modello è l'unico errore che viene trattato in maniera rigorosa nel processo di identificazione
    \item Per identificazione si intende l'assegnazione dei valori ai parametri incogniti del modello sulla base di conoscenze a priori 
    \item L'identificabilità a priori deve essere verificata sia nel modello input/output che nel modello strutturale 
    \item Il metodo della funzione di trasferimento permette di verificare l'identificabilità a priori dei modelli lineari 
    \item I singoli parametri del modello sono univocamente identificabili a priori se il riassunto esaustivo ammette un'unica soluzione 
    \item Un modello è non identificabile a priori se almeno uno dei suoi parametri non è identificabile   
    \item Per problema di stima si intende l'assegnazione di valori numerici ai parametri incogniti del modello a partire dalle misure effettuate sul sistema 
    \item Per stimare i parametri si utilizza sempre la regressione lineare sia nel caso di modello lineare nei parametri, che non lineare 
    \item Per residuo si intende la differenza tra l'ingresso e l'uscita del modello 
    \item La regressione lineare si basa sulla minimizzazione della funzione RSS o WRSS rispetto ai parametri incogniti 
    \item Il test dei residui permette di verificare la correttezza del modello impiegato
\end{enumerate}
\end{multicols}
